\section{Array Sections}
\label{sec:array_sections_syntax}
\index{array sections!syntax}

An array section designates a subset of the elements in an array
and may be used in OpenMP directives.
Array sections may be specified on multidimensional arrays.

\begin{ccppspecific}
For C/C++ the syntax of an array section is:

\begin{indentedcodelist}
[\plc{ [lower-bound] }:\plc{ [length] [}:\plc{ stride] }]
\end{indentedcodelist}

When evaluated, an array section represents a set of integer values
as follows:

\{ \plc{lower-bound}, \plc{lower-bound} + \plc{stride}, \plc{lower-bound} + 2 * \plc{stride},... , \plc{lower-bound} + ((\plc{length} - 1) * \plc{stride}) \}

When the \plc{lower-bound} is absent it defaults to 0.
When the \plc{stride} is absent it defaults to 1.
When the \plc{length} is absent and the size of the dimension is known, it
defaults to $(\plc{size} - \plc{lower-bound})/\plc{stride}$,
where \plc{size} is the size of the array dimension.
When the \plc{length} is absent and the size of the dimension is unknown, 
the array section is an assumed-size array.

The following are examples of array sections:

\begin{boxedcode}
a[0:6]
a[0:6:1]
a[1:10]
a[1:]
a[:10:2]
b[10][:][:]
b[10][:][:0]
c[42][0:6][:]
c[42][0:6:2][:]
c[1:10][42][0:6]
S.c[:100]
p->y[:10]
this->a[:N]
(p+10)[:N]
\end{boxedcode}

Assume \ucode{a} is declared to be a 1-dimensional array with dimension
size 11. The first two examples are equivalent, and the third and fourth
examples are equivalent. The fifth example specifies a stride of 2 and
therefore is not contiguous.

Assume \ucode{b} is declared to be a pointer to a 2-dimensional array with
dimension sizes 10 and 10. The sixth example refers to all elements of the
2-dimensional array given by \ucode{b[10]}. The seventh example is an
array section with zero length.

Assume \ucode{c} is declared to be a 3-dimensional array with dimension
sizes 50, 50, and 50.  The eighth example is contiguous, while the ninth
and tenth examples are not contiguous.

The final four examples show array sections that are formed from
more general array bases.

The following are examples that are non-conforming array sections:

\begin{boxedcode}
s[:10].x
p[:10]->y
*(xp[:10])
\end{boxedcode}

For all three examples, a base language operator is applied in an
undefined manner to an array section. The only operator that may be
applied to an array section is a subscript operator for which the
array section appears as the postfix expression.
\end{ccppspecific}

\begin{fortranspecific}
Fortran has built-in support for array sections with some
restrictions applied to their use in OpenMP directives.
For example, if a stride is specified, it must be positive.
See the OpenMP Specification for details.
\end{fortranspecific}
